% !TeX root = ../main.tex

\section*{Homework \#11 [2025-11-18]}

\begin{problem}[label = prob.11.1]\leavevmode
  \begin{enumext}
    \item Write the time-ordered product
    \[
      \mathcal T_\epsilon \{a_{k\sigma}(t_1)a_{l\sigma'}^\dagger(t_2)
                            a_{m\sigma}(t_3) a_{n\sigma'}^\dagger(t_3)\}
    \]
    in terms of normal products and suitable contractions.
    \item Express the expectation value of the time-ordered product in part
    (a) in the ground state $\ket|\eta_0>$ of the unperturbed system in terms
    of products of the free causal Green's functions.
  \end{enumext}
\end{problem}
\begin{solution}\leavevmode
  \begin{enumext}
    \item We define the time-ordered operators
    \[
      A = a_{k\sigma}(t_1), \quad
      B^\dagger = a_{l\sigma'}^\dagger(t_2), \quad
      C = a_{m\sigma}(t_3),\quad
      D^\dagger = a_{n\sigma'}^\dagger(t_3)
    \]
    Using Wick's theorem for fermions, the time-ordered product can be
    written as
    \begin{multline*}
      \mathcal T\{AB^\dagger CD^\dagger\}
  = \mathcal N\{AB^\dagger CD^\dagger\} +
        \sum_\text{Single contractions}
        (-1)^\sigma (\text{contraction}) \\
        \times (\text{Normal ordered remainder})
      + \sum_\text{Double contractions} (-1)^\sigma
        (\text{Product of contractions})
    \end{multline*}
    where the contraction between two operators, $\hat c_1$ and $\hat c_2$,
    for example, is defined as the expectation value of the time-ordered product
    in reference vacuum
    \[
      \wick{\c1{\hat c}_1\c1{\hat c}_2}
        \equiv \braket<T_\epsilon\{\hat c_1\hat c_2\}>_0.
    \]
    Note that the contractions $\wick{\c1A\c1C}$
    and $\wick{\c1B^\dagger\c1D^\dagger}$ vanish
    since they are all annihilation operators or creation operators.
    The non-zero contractions are listed as follows
    \begin{enumext}
      \item Single contraction $(1,2)$:
      $(1,2,3,4) \to (1,2,3,4)$ contains two swaps: the sign is $(-1)^0 = +$.
      \item Single contraction $(3,4)$:
      $(1,2,3,4) \to (3,4,1,2)$ contains two swaps: the sign is $(-1)^4 = +$.
      \item Single contraction $(1,4)$:
      $(1,2,3,4) \to (1,4,2,3)$ contains two swaps: the sign is $(-1)^2 = +$.
      \item Single contraction $(2,3)$:
      $(1,2,3,4) \to (2,3,1,4)$ contains two swaps: the sign is $(-1)^2 = +$.
      \item Double contraction $(1,2)$ and $(3,4)$:
      the sign is $(-1)^{0+4} = +$.
      \item Double contraction $(1,4)$ and $(2,3)$:
      the sign is $(-1)^{2+2} = +$.
    \end{enumext}
    where $A$, $B$, $C$, $D$ have already defined before.
    Substitute the contraction into the Wick's theorem
    \begin{multline*}
      \mathcal T_\epsilon\{AB^\dagger CD^\dagger\}
    = :AB^\dagger CD^\dagger:
    + \wick{\c1A\c1B^\dagger}:CD^\dagger:
    + \wick{\c1C\c1D^\dagger} :AB^\dagger:\\
    + \wick{\c1A\c1D^\dagger} :B^\dagger C:
    + \wick{\c1B^\dagger\c1C} :AD^\dagger:
    + \wick{\c1A\c1B^\dagger\c1C\c1D^\dagger}
    + \wick{\c1A\c1D^\dagger\c1B^\dagger\c1C}
    \end{multline*}
    Restoring the explicit indices
    \begin{align*}
      \mathcal T_\epsilon \{a_{k\sigma}(t_1) a_{l\sigma'}^\dagger(t_2)
                            a_{m\sigma}(t_3) a_{n\sigma'}^\dagger(t_3)\}
  = &:a_{k\sigma}(t_1) a_{l\sigma'}^\dagger(t_2)
      a_{m\sigma}(t_3) a_{n\sigma'}^\dagger(t_3):\\
    & + \braket<\mathcal T_\epsilon\{
          a_{k\sigma}(t_1) a_{l\sigma'}^\dagger(t_2)\}>_0
         :a_{m\sigma(t_3)} a_{n\sigma'}^\dagger(t_3):\\
    & + \braket<\mathcal T_\epsilon\{
          a_{m\sigma(t_3)} a_{n\sigma'}^\dagger(t_3)\}>_0
         :a_{k\sigma}(t_1) a_{l\sigma'}^\dagger(t_2):\\
    & + \braket<\mathcal T_\epsilon\{
          a_{k\sigma}(t_1) a_{n\sigma'}^\dagger(t_3)\}>_0
         :a_{l\sigma'}^\dagger(t_2) a_{m\sigma(t_3)}:\\
    & + \braket<\mathcal T_\epsilon\{
          a_{l\sigma'}^\dagger(t_2)a_{m\sigma}(t_3)\}>_0
         :a_{k\sigma}(t_1) a_{n\sigma'}^\dagger(t_3):\\
    & + \braket<\mathcal T_\epsilon\{
          a_{k\sigma}(t_1) a_{l\sigma'}^\dagger(t_2)\}>_0
        \braket<\mathcal T_\epsilon\{
          a_{m\sigma}(t_3) a_{n\sigma'}^\dagger(t_3)\}>_0\\
    & + \braket<\mathcal T_\epsilon\{
          a_{k\sigma}(t_1) a_{n\sigma'}^\dagger(t_3)\}>_0
        \braket<\mathcal T_\epsilon\{
          a_{l\sigma'}^\dagger(t_2) a_{m\sigma}(t_3)\}>_0
    \end{align*}
    \item The expectation value in the ground state $\ket|\eta_0>$ is
    contributed by the \emph{fully contracted terms}, i.e., the last two double
    contraction terms. Using the definition of the free causal Green's function
    \[
      \wick{\c1 a_{k\sigma}(t)\c1 a_{l\sigma'}^\dagger(t')}
    = \iu G_{k\sigma,l\sigma'}(t,t')
    = \braket<\eta_0|\mathcal T_\epsilon
        \{a_{k\sigma}(t) a_{l\sigma}^\dagger(t')\}|\eta_0>
    \]
    To uniform the order of the annihilation and creation operators, using
    the ``anti-commutation relation''
    \[
      \wick{\c1 a_{l\sigma'}^\dagger(t_2) \c1a_{m\sigma}(t_3)}
    = -\iu G_{m\sigma,l\sigma'}(t_3,t_2).
    \]
    Therefore, the expectation value
    \begin{align*}
      \braket<\eta_0|\mathcal T_\epsilon\{AB^\dagger CD^\dagger\}|\eta_0> &
    = [\iu G_{k\sigma,l\sigma'}^0(t_1,t_2)]
      [\iu G_{m\sigma,n\sigma'}^0(t_3,t_3)] +
      [\iu G_{k\sigma,n\sigma'}^0(t_1,t_3)]
      [-\iu G_{m\sigma,l\sigma'}^0(t_3,t_2)]\\
  & = -G_{k\sigma,l\sigma'}^0(t_1,t_2)
      G_{m\sigma,n\sigma'}^0(t_3,t_3)
    + G_{k\sigma,n\sigma'}^0(t_1,t_3)
      G_{m\sigma,l\sigma'}^0(t_3,t_2).
    \end{align*}
  \end{enumext}
\end{solution}

\begin{problem}
  Evaluate explicitly the expectation value of the time-ordered product in
  part 1 of \Autoref{prob.11.1} in the ground state $\eta\ket|\eta_0>$ of the
  unperturbed system for the special case $k = l = m = n$, $\sigma = \sigma'$
  for
  \begin{enumext}[columns = 2]
    \item $t_1 > t_2 > t_3$,
    \item $t_1 > t_3 > t_2$.
  \end{enumext}
  Check the results by direct calculation of the expectation values, i.e.,
  without using Wick's theorem.
\end{problem}
\begin{solution}
  Since $k = l = m = n$, $\sigma = \sigma'$, so the labels of the fermion
  operators make no sense here. Calculate an identity first. Consider a very
  common situation: The phase factors of the fermion operators are
  time-dependent, i.e.,
  \[
    a(t) = \upe^{-\iu\int_{-\infty}^t \d t'\epsilon(t')} a(0), \qq{and}
    a^\dagger(t') = \upe^{\iu\int_{-\infty}^t \d t'\epsilon(t')} a^\dagger(0).
  \]
  Denote $f(\epsilon^0) = \braket<a(0) a^\dagger(0)>$,
  we can express the double time expectation
  \[
    \braket<c(t) c^\dagger(t')> = \braket<c^\dagger(t') c(t)>
  = f(\epsilon^0)\upe^{-\iu\int_{t'}^t \d t_1 \epsilon(t_1)}
  \]
  The following calculations mainly use the anti-commutation relations
  \[
    aa^\dagger = 1 = a^\dagger a = 1 - \hat n, \quad
    a^\dagger aa^\dagger a = a^\dagger a
  \]
  \begin{enumext}
    \item Expand the time-order operators first, then substitute the identity
    \begin{align*}
      \braket<\eta_0|\mathcal T_\epsilon\{AB^\dagger CD^\dagger\}|\eta_0>
  & = \braket<a(t_1) a^\dagger(t_2)>
      \braket<a(t_3) a^\dagger(t_3)>
    + \braket<a(t_1) a^\dagger(t_3)>
     \braket<a^\dagger(t_2) a(t_3)>\\
  & = (1 - n)^2
      \upe^{-\iu\int_{t_2}^{t_1} \d t'\epsilon(t')}
      \upe^{-\iu\int_{t_3}^{t_3} \d t'\epsilon(t')}
    + (1 - n) \cdot n
      \upe^{-\iu\int_{t_3}^{t_1} \d t'\epsilon(t')}
      \upe^{-\iu\int_{t_2}^{t_3} \d t'\epsilon(t')}\\
  & = (1 - n) \upe^{-\iu\int_{t_2}^{t_1} \d t'\epsilon(t')}
      \xlongequal{\pdv\epsilon/t = 0} (1 - n)
      \upe^{-\iu\epsilon(t_1 - t_2)}
    \end{align*}
    By direct calculation, starting from the original expression. Obviously,
    they have the same phase factor
    \begin{align*}
      \braket<
        \mathcal T_\epsilon\{a(t_1) a^\dagger(t_2) a(t_3) a^\dagger(t_3)\}>_0
      \big|_{t_1>t_2>t_3}
  & = \braket<a a^\dagger a a^\dagger>
      \upe^{-\iu\int_{t_2}^{t_1} \d t'\epsilon(t')}
    = \braket<(1 - a^\dagger a)(1 - a^\dagger a)>
      \upe^{-\iu\int_{t_2}^{t_1} \d t'\epsilon(t')}\\
  & = (1 - 2aa^\dagger + a^\dagger a)
      \upe^{-\iu\int_{t_2}^{t_1} \d t'\epsilon(t')}
    = (n - 1) \upe^{-\iu\int_{t_2}^{t_1} \d t'\epsilon(t')}
    \end{align*}
    The result matches.
    \item Different from (a), the last term will have an extra minus sign
    \begin{multline*}
      \braket<\eta_0|\mathcal T_\epsilon\{AB^\dagger CD^\dagger\}|\eta_0>
    = \braket<a(t_1) a^\dagger(t_2)>
      \braket<a(t_3) a^\dagger(t_3)>
    + \braket<a(t_1) a^\dagger(t_3)>
      \braket<-a(t_3) a^\dagger(t_2)>\\
    = \upe^{-\iu\int_{t_2}^{t_1} \d t'\epsilon(t')}(
      \braket<aa^\dagger> \braket<aa^\dagger>
    - \braket<aa^\dagger> \braket<aa^\dagger>) = 0
    \end{multline*}
    In the direct calculation,
    $\braket<\mathcal T\{a(t_1)a^\dagger(t_2)a(t_3)a^\dagger(t_3)\}>$
    gives $\braket<a(t_1)a(t_3)a^\dagger(t_3)a^\dagger(t_2)>$.
    The term $a(t_1)a(t_3)$ must make the whole expression vanish.
    The result matches.
  \end{enumext}
\end{solution}